\documentclass{article}
\usepackage{amsmath,amssymb,amsthm}
\usepackage{algorithm}
\usepackage{algpseudocode}
\usepackage{bm}
\usepackage{hyperref}
\newtheorem{theorem}{Theorem}
\newtheorem{lemma}{Lemma}
\newcommand{\E}{\mathbb{E}}
\newcommand{\R}{\mathbb{R}}

\begin{document}

\section*{Differentially Private Stochastic Gradient Descent (DP‑SGD) with Gaussian Noise}

\section*{Mathematical Formulation}
Let $f:\R^{d}\to\R$ be a differentiable, possibly non‑convex loss function.  
At iteration $t\in\{0,1,\dots,T-1\}$ we draw a stochastic gradient
\[
g_{t}\;=\;\nabla f(x_{t})\;+\;\xi_{t},
\]
where the noise $\xi_{t}\sim\mathcal{N}(0,\sigma^{2}I_{d})$ is independent across
iterations and $I_{d}$ denotes the $d\times d$ identity matrix.  

Given a (possibly time‑varying) stepsize $\eta_{t}>0$, the DP‑SGD update is
\begin{equation}
\boxed{%
x_{t+1}\;=\;x_{t}\;-\;\eta_{t}\,g_{t}
\;=\;x_{t}\;-\;\eta_{t}\bigl(\nabla f(x_{t})+\xi_{t}\bigr).}
\tag{1}
\end{equation}
The privacy budget $(\varepsilon,\delta)$ is tracked by a moments accountant
(or any standard composition theorem) using the known variance $\sigma^{2}$
and the sampling probability $q$ (if mini‑batches are used).  For the
purpose of convergence analysis we treat $\sigma^{2}$ as a fixed constant.

\section*{Pseudocode}
\begin{algorithm}[H]
\caption{DP‑SGD with Gaussian Noise}
\begin{algorithmic}[1]
\Require Initial point $x_{0}\in\R^{d}$, stepsizes $\{\eta_{t}\}_{t=0}^{T-1}$,
noise level $\sigma>0$, total iterations $T$.
\For{$t=0$ \textbf{to} $T-1$}
    \State Compute the (true) gradient $\nabla f(x_{t})$.
    \State Sample $\xi_{t}\sim\mathcal{N}(0,\sigma^{2}I_{d})$.
    \State Form noisy gradient $g_{t}\gets\nabla f(x_{t})+\xi_{t}$.
    \State Update $x_{t+1}\gets x_{t}-\eta_{t}g_{t}$.
    \State (Optional) Update privacy accountant with $\sigma$ and $q$.
\EndFor
\State \Return $x_{T}$ (or the average $\bar x_{T}= \tfrac1T\sum_{t=0}^{T-1}x_{t}$).
\end{algorithmic}
\end{algorithm}

\section*{Dry Run Example}
Consider the one‑dimensional quadratic
\[
f(w)=(w-3)^{2},\qquad \nabla f(w)=2(w-3).
\]
Set $w_{0}=0$, constant stepsize $\eta=0.1$, and fix a noise realization
$\xi_{0}=0.2$, $\xi_{1}=-0.1$, $\xi_{2}=0.05$ (all drawn from $\mathcal N(0,\sigma^{2})$).

\begin{center}
\begin{tabular}{c|c|c|c}
$t$ & $\nabla f(w_{t})$ & $g_{t}= \nabla f(w_{t})+\xi_{t}$ & $w_{t+1}=w_{t}-\eta g_{t}$ \\ \hline
0 & $2(0-3)=-6$ & $-6+0.2=-5.8$ & $0-0.1(-5.8)=0.58$ \\[4pt]
1 & $2(0.58-3)=-4.84$ & $-4.84-0.1=-4.94$ & $0.58-0.1(-4.94)=1.074$ \\[4pt]
2 & $2(1.074-3)=-3.852$ & $-3.852+0.05=-3.802$ & $1.074-0.1(-3.802)=1.4542$
\end{tabular}
\end{center}
Corresponding objective values:
$f(w_{0})=9$, $f(w_{1})\approx6.17$, $f(w_{2})\approx3.38$, $f(w_{3})\approx2.12$.
The noisy updates push the iterate towards the minimiser $w^{*}=3$ while
introducing stochastic perturbations.

\newpage
\section*{Assumptions}
\begin{enumerate}
\item \textbf{Smoothness.} $f$ is $L$‑smooth, i.e.
\[
\forall x,y\in\R^{d}:\quad
f(y)\le f(x)+\langle\nabla f(x),y-x\rangle+\frac{L}{2}\|y-x\|^{2}.
\tag{A1}
\]
\item \textbf{Bounded Below.} $f^{*}\triangleq\inf_{x}f(x)>-\infty$.
\item \textbf{Gradient Variance.} The Gaussian noise satisfies
\[
\E[\xi_{t}]=0,\qquad \E[\|\xi_{t}\|^{2}]=d\sigma^{2}.
\tag{A2}
\]
\item \textbf{Stepsize.} The stepsizes are constant $\eta_{t}\equiv\eta$ with
$0<\eta\le\frac{1}{L}$.
\end{enumerate}

\section*{Main Convergence Theorem}
\begin{theorem}[Expected Gradient Norm for Non‑Convex DP‑SGD]\label{thm:main}
Under assumptions (A1)–(A4) and after $T$ iterations of \eqref{1},
the averaged iterate $\bar x_{T}\triangleq\frac{1}{T}\sum_{t=0}^{T-1}x_{t}$ satisfies
\[
\frac{1}{T}\sum_{t=0}^{T-1}\E\!\bigl[\|\nabla f(x_{t})\|^{2}\bigr]\;
\le\;
\underbrace{\frac{2\bigl(f(x_{0})-f^{*}\bigr)}{\eta T}}_{\text{deterministic decay}}
\;+\;
\underbrace{L\eta\,d\sigma^{2}}_{\text{privacy‑induced variance}}.
\tag{2}
\]
Consequently, choosing $\eta=\Theta\!\bigl(1/\sqrt{T}\bigr)$ yields
\[
\min_{0\le t<T}\E\!\bigl[\|\nabla f(x_{t})\|^{2}\bigr]
=O\!\left(\frac{1}{\sqrt{T}}\right)+O\!\left(\frac{d\sigma^{2}}{\sqrt{T}}\right).
\]
\end{theorem}

\section*{Theoretical Properties}
\begin{itemize}
\item \textbf{Stability.} The bound \eqref{2} shows that as long as
$\eta\le 1/L$, the iterates remain stable; the term $L\eta d\sigma^{2}$
captures the steady‑state noise floor caused by differential privacy.
\item \textbf{Hyperparameter Sensitivity.} Smaller $\eta$ reduces the
privacy‑induced variance but slows deterministic decay; the optimal
choice balances the two terms, giving $\eta^{\star}= \sqrt{\frac{2(f(x_{0})-f^{*})}{L d\sigma^{2}T}}$.
\item \textbf{Comparison.} When $\sigma=0$ the bound collapses to the
standard non‑convex SGD result $\frac{2(f(x_{0})-f^{*})}{\eta T}$.
With Gaussian noise the additional $L\eta d\sigma^{2}$ term is
unavoidable under $(\varepsilon,\delta)$‑DP.
\end{itemize}

\section*{Discussion}
The theorem establishes that DP‑SGD converges to a stationary point
in expectation at the same $O(1/\sqrt{T})$ rate as vanilla SGD, up to an
additive term proportional to the noise variance.  This term is the
price paid for privacy and can be mitigated by increasing the batch
size (which reduces the effective $\sigma$ after accounting for sampling)
or by decaying the stepsize.

\newpage
\section*{Preliminaries (Definitions and Lemmas)}
\begin{lemma}[Smoothness Inequality]\label{lem:smooth}
If $f$ satisfies (A1), then for any $x\in\R^{d}$ and any vector $u$,
\[
f(x-u)\le f(x)-\langle\nabla f(x),u\rangle+\frac{L}{2}\|u\|^{2}.
\tag{3}
\]
\end{lemma}
\begin{proof}
Apply (A1) with $y=x-u$ and rearrange. \hfill $\square$
\end{proof}

\begin{lemma}[Noise Moment]\label{lem:noise}
For $\xi\sim\mathcal N(0,\sigma^{2}I_{d})$,
\[
\E[\xi]=0,\qquad \E[\|\xi\|^{2}]=d\sigma^{2}.
\tag{4}
\]
\end{lemma}
\begin{proof}
Standard properties of the multivariate normal distribution. \hfill $\square$
\end{proof}

\section*{Main Proof (Step‑by‑Step Derivation)}
We prove Theorem~\ref{thm:main}.  Throughout the proof expectations are
taken w.r.t.\ all randomness up to the current iteration.

\noindent\textbf{[Step 1] Apply smoothness to the update.}  
Using Lemma~\ref{lem:smooth} with $u=\eta g_{t}$ and $x=x_{t}$,
\[
f(x_{t+1})\le f(x_{t})-\eta\langle\nabla f(x_{t}),g_{t}\rangle
+\frac{L\eta^{2}}{2}\|g_{t}\|^{2}.
\tag{5}
\]

\noindent\textbf{[Step 2] Expand $g_{t}$.}  
Recall $g_{t}= \nabla f(x_{t})+\xi_{t}$, thus
\[
\langle\nabla f(x_{t}),g_{t}\rangle
= \|\nabla f(x_{t})\|^{2}+\langle\nabla f(x_{t}),\xi_{t}\rangle,
\qquad
\|g_{t}\|^{2}= \|\nabla f(x_{t})\|^{2}
+2\langle\nabla f(x_{t}),\xi_{t}\rangle+\|\xi_{t}\|^{2}.
\tag{6}
\]

\noindent\textbf{[Step 3] Substitute (6) into (5).}
\[
\begin{aligned}
f(x_{t+1})
&\le f(x_{t})
-\eta\bigl(\|\nabla f(x_{t})\|^{2}+\langle\nabla f(x_{t}),\xi_{t}\rangle\bigr) \\
&\qquad +\frac{L\eta^{2}}{2}
\bigl(\|\nabla f(x_{t})\|^{2}+2\langle\nabla f(x_{t}),\xi_{t}\rangle+\|\xi_{t}\|^{2}\bigr).
\end{aligned}
\tag{7}
\]

\noindent\textbf{[Step 4] Collect terms.}
\[
\begin{aligned}
f(x_{t+1})
&\le f(x_{t})
-\Bigl(\eta-\frac{L\eta^{2}}{2}\Bigr)\|\nabla f(x_{t})\|^{2}\\
&\quad -\Bigl(\eta-\!L\eta^{2}\Bigr)\langle\nabla f(x_{t}),\xi_{t}\rangle
+\frac{L\eta^{2}}{2}\|\xi_{t}\|^{2}.
\end{aligned}
\tag{8}
\]

\noindent\textbf{[Step 5] Take expectation conditioned on $x_{t}$.}  
Because $\xi_{t}$ is independent of $x_{t}$ and $\E[\xi_{t}]=0$,
\[
\E\!\bigl[\langle\nabla f(x_{t}),\xi_{t}\rangle\mid x_{t}\bigr]=0.
\tag{9}
\]
Moreover, by Lemma~\ref{lem:noise},
\[
\E\!\bigl[\|\xi_{t}\|^{2}\mid x_{t}\bigr]=d\sigma^{2}.
\tag{10}
\]

\noindent\textbf{[Step 6] Apply (9)–(10) to (8).}
\[
\E\!\bigl[f(x_{t+1})\mid x_{t}\bigr]
\le f(x_{t})-\Bigl(\eta-\frac{L\eta^{2}}{2}\Bigr)\|\nabla f(x_{t})\|^{2}
+\frac{L\eta^{2}}{2}\,d\sigma^{2}.
\tag{11}
\]

\noindent\textbf{[Step 7] Unconditional expectation.}  
Taking full expectation and using the law of total expectation,
\[
\E\!\bigl[f(x_{t+1})\bigr]
\le \E\!\bigl[f(x_{t})\bigr]
-\Bigl(\eta-\frac{L\eta^{2}}{2}\Bigr)\E\!\bigl[\|\nabla f(x_{t})\|^{2}\bigr]
+\frac{L\eta^{2}}{2}\,d\sigma^{2}.
\tag{12}
\]

\noindent\textbf{[Step 8] Simplify coefficient using $\eta\le 1/L$.}  
Since $0<\eta\le 1/L$, we have $L\eta\le1$ and thus
\[
\eta-\frac{L\eta^{2}}{2}\;\ge\;\frac{\eta}{2}.
\tag{13}
\]

\noindent\textbf{[Step 9] Plug (13) into (12).}
\[
\E\!\bigl[f(x_{t+1})\bigr]
\le \E\!\bigl[f(x_{t})\bigr]
-\frac{\eta}{2}\E\!\bigl[\|\nabla f(x_{t})\|^{2}\bigr]
+\frac{L\eta^{2}}{2}\,d\sigma^{2}.
\tag{14}
\]

\noindent\textbf{[Step 10] Telescoping sum.}  
Sum \eqref{14} from $t=0$ to $T-1$:
\[
\E\!\bigl[f(x_{T})\bigr]
\le f(x_{0})
-\frac{\eta}{2}\sum_{t=0}^{T-1}\E\!\bigl[\|\nabla f(x_{t})\|^{2}\bigr]
+\frac{L\eta^{2}}{2}\,Td\sigma^{2}.
\tag{15}
\]

\noindent\textbf{[Step 11] Lower bound $f(x_{T})\ge f^{*}$.}
\[
f^{*}\le \E\!\bigl[f(x_{T})\bigr].
\tag{16}
\]

\noindent\textbf{[Step 12] Rearrange (15)–(16).}
\[
\frac{\eta}{2}\sum_{t=0}^{T-1}\E\!\bigl[\|\nabla f(x_{t})\|^{2}\bigr]
\le f(x_{0})-f^{*}
+\frac{L\eta^{2}}{2}\,Td\sigma^{2}.
\tag{17}
\]

\noindent\textbf{[Step 13] Divide by $\eta T/2$.}
\[
\frac{1}{T}\sum_{t=0}^{T-1}\E\!\bigl[\|\nabla f(x_{t})\|^{2}\bigr]
\le \frac{2\bigl(f(x_{0})-f^{*}\bigr)}{\eta T}
+L\eta\,d\sigma^{2}.
\tag{18}
\]
Equation \eqref{18} is exactly the bound \eqref{2} claimed in Theorem~\ref{thm:main}.
\hfill $\square$

\section*{Rate Derivation (With Constants)}
From \eqref{18}, set $\eta = \min\!\bigl\{1/L,\sqrt{\frac{2(f(x_{0})-f^{*})}{L d\sigma^{2}T}}\bigr\}$.
If $T$ is large enough so that the square‑root term is $\le 1/L$, then
\[
\frac{1}{T}\sum_{t=0}^{T-1}\E\!\bigl[\|\nabla f(x_{t})\|^{2}\bigr]
\le
O\!\Bigl(\frac{1}{\sqrt{T}}\Bigr)+O\!\Bigl(\frac{d\sigma^{2}}{\sqrt{T}}\Bigr).
\]
Hence the expected squared gradient norm decays as $O(T^{-1/2})$.

\section*{Special Cases}
\begin{enumerate}
\item \textbf{No privacy noise ($\sigma=0$).} The bound reduces to the classic
non‑convex SGD result
\[
\frac{1}{T}\sum_{t=0}^{T-1}\E[\|\nabla f(x_{t})\|^{2}]
\le \frac{2(f(x_{0})-f^{*})}{\eta T}.
\]
\item \textbf{Constant stepsize $\eta=1/L$.} Plugging into \eqref{18} yields
\[
\frac{1}{T}\sum_{t=0}^{T-1}\E[\|\nabla f(x_{t})\|^{2}]
\le \frac{2L\bigl(f(x_{0})-f^{*}\bigr)}{T}
+ d\sigma^{2}.
\]
Thus the algorithm converges to a neighbourhood of a stationary point,
the radius of which is $O(d\sigma^{2})$.
\item \textbf{Mini‑batch sampling.} If each iteration uses a batch of size $b$
with sampling probability $q=b/n$, the effective noise variance becomes
$\sigma^{2}_{\text{eff}}= \sigma^{2}/b$, improving the bound by a factor $1/b$.
\end{enumerate}

\end{document}